\chapter{Im và Nói}
\markboth{Im và Nói}{Silence and Speech}

\section{Im lặng đúng lúc để không làm đau thêm.}
\begin{truyen}{Im lặng đúng lúc để không làm đau thêm.}{Silence can be a form of care.}
\chuhoa{K}{hi người thân đang quá nóng, cô biết mọi lời giải thích lúc đó chỉ làm lửa bùng cao hơn. Cô chọn im, không phải vì hết lý lẽ, mà vì còn muốn giữ nhau. Có sự im lặng là trốn tránh, nhưng cũng có sự im lặng là lòng thương biết chờ thời điểm đúng.}
\textit{(When a loved one was too angry, she knew that any explanation at that moment would only raise the flames. She chose silence, not because she lacked reasons, but because she still wanted to preserve the relationship. Some silence is avoidance, but some silence is care waiting for the right time.)}
\end{truyen}
\begin{baihoc}
Không phải lúc nào nói ra ngay cũng là can đảm; đôi khi can đảm là biết giữ miệng để giữ người.
\textit{(Speaking immediately is not always courage; sometimes courage is guarding your mouth to protect the relationship.)}
\end{baihoc}
\begin{ghinhoanh}
\textbf{Short English to remember:}

Silence can protect love.
\end{ghinhoanh}
\begin{tuhocnhanh}
1. Em có tình huống nào đáng lẽ nên nói ít lại? \textit{(Is there a situation where you should have said less?)}

2. Với em, im lặng nào là khôn chứ không phải hèn? \textit{(For you, what kind of silence is wise rather than cowardly?)}
\end{tuhocnhanh}
\ngancach

\section{Nói đúng lúc để ngăn một điều sai tiếp diễn.}
\begin{truyen}{Nói đúng lúc để ngăn một điều sai tiếp diễn.}{Silence is not always kindness.}
\chuhoa{T}{hấy bạn bị trêu chọc quá mức, em không cười theo như mọi lần. Em lên tiếng ngắn gọn rằng chuyện đó nên dừng lại. Có những lúc im lặng không còn là hiền nữa, mà thành đồng lõa.}
\textit{(Seeing a classmate mocked too far, the student did not laugh along as usual. Instead, there was a short statement that it needed to stop. There are moments when silence is no longer gentleness, but complicity.)}
\end{truyen}
\begin{baihoc}
Biết nói là một trách nhiệm khi điều sai đang diễn ra trước mặt.
\textit{(Knowing when to speak is a responsibility when wrong is happening in front of you.)}
\end{baihoc}
\begin{ghinhoanh}
\textbf{Short English to remember:}

Silence can become consent.
\end{ghinhoanh}
\begin{tuhocnhanh}
1. Em đã từng im lặng ở một chỗ đáng lẽ nên lên tiếng chưa? \textit{(Have you ever stayed silent where you should have spoken?)}

2. Điều gì làm em sợ nói ra điều đúng? \textit{(What makes you afraid to speak what is right?)}
\end{tuhocnhanh}
\ngancach

\section{Nói nhiều để che nỗi trống bên trong.}
\begin{truyen}{Nói nhiều để che nỗi trống bên trong.}{Too many words can hide a hollow center.}
\chuhoa{C}{ó người lúc nào cũng kể, cũng đùa, cũng chen lời, như thể sợ một khoảng lặng nào đó. Chỉ khi ở một mình, họ mới thấy trong mình có quá nhiều câu hỏi chưa dám đối diện. Đôi khi cái ồn ào của miệng chỉ là cách chạy khỏi sự thật của lòng.}
\textit{(Some people are always talking, joking, interrupting, as if afraid of silence itself. Only when alone do they notice the many questions inside that remain unfaced. Sometimes the noise of the mouth is simply a way of fleeing the truth of the heart.)}
\end{truyen}
\begin{baihoc}
Không phải nói nhiều là giao tiếp tốt; nhiều khi đó chỉ là né tránh nội tâm.
\textit{(Talking a lot is not always good communication; often it is avoidance of the inner world.)}
\end{baihoc}
\begin{ghinhoanh}
\textbf{Short English to remember:}

Noise can hide emptiness.
\end{ghinhoanh}
\begin{tuhocnhanh}
1. Em có dùng sự bận rộn hay nói nhiều để né điều gì không? \textit{(Do you use busyness or excessive talking to avoid something?)}

2. Nếu im xuống một chút, điều gì trong em sẽ hiện ra? \textit{(If you became quieter, what inside you would appear?)}
\end{tuhocnhanh}
\ngancach

\section{Im để nghe người khác thật sự.}
\begin{truyen}{Im để nghe người khác thật sự.}{Listening begins where self-display stops.}
\chuhoa{M}{ột thầy giáo khi nghe học sinh kể chuyện gia đình không chen vào giảng ngay. Thầy để em nói hết, cả những chỗ ngập ngừng lẫn những chỗ lộn xộn. Chính sự im lặng biết lắng nghe ấy làm đứa học trò cảm thấy mình được tôn trọng.}
\textit{(When a student shared family pain, a teacher did not interrupt with lessons too quickly. He allowed the full story, including the hesitant and messy parts. That listening silence made the student feel respected.)}
\end{truyen}
\begin{baihoc}
Im đúng cách mở ra không gian để người khác được là chính họ.
\textit{(Silence used well opens space for others to be themselves.)}
\end{baihoc}
\begin{ghinhoanh}
\textbf{Short English to remember:}

Listening needs silence.
\end{ghinhoanh}
\begin{tuhocnhanh}
1. Khi người khác nói, em nghe để hiểu hay nghe để chờ phản hồi? \textit{(When others speak, do you listen to understand or only to reply?)}

2. Em có thể tập im hơn ở cuộc trò chuyện nào? \textit{(In which conversation could you practice more silence?)}
\end{tuhocnhanh}
\ngancach

\section{Nói một câu thật thay cho nhiều lời xã giao.}
\begin{truyen}{Nói một câu thật thay cho nhiều lời xã giao.}{One honest sentence can outweigh many polite words.}
\chuhoa{K}{hi bạn sắp bỏ cuộc, em không nói những câu động viên rập khuôn. Em chỉ nói: “Tớ biết cậu đang mệt thật, nhưng tớ sẽ ngồi với cậu đến hết phần này.” Một câu thật, có trách nhiệm, đôi khi cứu người hơn cả chục câu hay mà rỗng.}
\textit{(When a friend was about to give up, the student did not use formulaic encouragement. Instead came one sentence: “I know you are really tired, but I will stay with you through this part.” One honest, responsible sentence can help more than ten beautiful but empty ones.)}
\end{truyen}
\begin{baihoc}
Lời nói có giá trị không nằm ở độ hoa mỹ, mà ở độ thật và độ chịu trách nhiệm.
\textit{(The value of speech lies not in elegance, but in honesty and responsibility.)}
\end{baihoc}
\begin{ghinhoanh}
\textbf{Short English to remember:}

Honest words carry weight.
\end{ghinhoanh}
\begin{tuhocnhanh}
1. Em có đang nói nhiều câu đẹp nhưng ít câu thật không? \textit{(Are you speaking many nice words but too few honest ones?)}

2. Câu nào em cần nói thật với một người quan trọng? \textit{(What honest sentence do you need to say to someone important?)}
\end{tuhocnhanh}
\ngancach
